\section{Example Exam Questions}

\noindent The following is a non-exhaustive list of possible oral exam questions, organized by lecture topic, as given in the AE4262 course guidelines (April 17, 2026).

\subsection{Thermodynamics and Kinetics}

\subsubsection{Explain different types of combustion.}

\subsubsection{Explain the different types of enthalpies and how to determine the enthalpy of a reaction.}

\subsubsection{What is heat capacity and why does it vary?}

\subsubsection{What is adiabatic flame temperature and how does it vary with parameters like equivalence ratio, inlet temperature, pressure, fuel type, etc.?}

\subsubsection{What is Gibbs energy and how does it define spontaneity?}

\subsection{Laminar Transport}

\subsubsection{In steady state combustion studies, enthalpy $h$ is often chosen as the energy variable. What is the relation between enthalpy and temperature? Explain the physical meaning of the different terms in the equation for enthalpy (equation will be given). Explain the physical meaning of the different terms in the equation for temperature (equation will be given).}

\subsubsection{Transport of momentum due to viscosity is analogous to transport of heat by conduction and is analogous to transport of species by diffusion. Explain this and introduce relevant dimensionless numbers.}

\subsection{Turbulent Transport}

\subsubsection{What are relevant physical time scales in turbulent flow? How do they relate to chemical time scales?}

\subsubsection{Give a short description of the three main approaches to the computation of turbulent flames and specify their advantages and drawbacks.}

\subsubsection{Explain the main closure problems that arise when averaging the transport equations for an exothermic reacting flow.}

\subsubsection{Density-weighted averaging, also called Favre averaging, is an alternative way of averaging leading to a specific form of the averaged transport equations. What is Favre averaging? What is this specific form? Can it be expected to be a good approach for premixed flames / for non-premixed flames?}

\subsubsection{Explain what is meant by Large Eddy Simulation of turbulent flow. Also discuss the need for subgrid models of combustion processes in the context of LES.}

\subsection{Premixed Combustion}

\subsubsection{Discuss the structure of a flat laminar premixed flame.}

\subsubsection{Discuss the classical regime diagrams of turbulent premixed combustion (Borghi and Peters diagrams).}

\subsubsection{What is laminar flame speed? What is turbulent flame speed? How are they related?}

\subsubsection{What is turbulent flame speed closure?}

\subsubsection{How can a laminar premixed flame be described in an intrinsic way using a so-called progress variable? How can a turbulent premixed flame be described using a progress variable?}

\subsection{Non-Premixed Combustion}

\subsubsection{What is the typical structure of a laminar diffusion flame?}

\subsubsection{Why is there no such thing as a laminar flame speed of a diffusion flame?}

\subsubsection{Define mixture fraction. Which conditions have to be satisfied in order that mixture fraction becomes a useful variable? Explain the idea that, using mixture fraction, a non-premixed flame problem can be split into a mixing problem and a chemical problem.}

\subsubsection{Define mixture fraction. In which circumstances can enthalpy be considered equivalent to mixture fraction?}

\subsubsection{Explain how the influence of turbulent fluctuations of mixture fraction can be taken into account by making use of an assumed shape of its probability density function.}

\subsubsection{Explain why one is interested in calculating the variance of mixture fraction. How does the scalar dissipation rate influence the variance of mixture fraction?}

\subsubsection{Explain the steady laminar flamelet model of turbulent non-premixed combustion.}

\subsection{Combustion Modelling}

\subsubsection{Give a short description of the three main approaches to the computation of turbulent flames and specify their advantages and drawbacks.}

\subsubsection{Explain the main closure problems that arise when averaging the transport equations for an exothermic reacting flow.}

\subsubsection{Density-weighted averaging, also called Favre averaging, is an alternative way of averaging leading to a specific form of the averaged transport equations. What is Favre averaging? What is this specific form?}

\subsubsection{Explain what is meant by Large Eddy Simulation of turbulent flow. Also discuss the need for subgrid models of combustion processes in the context of LES.}

\subsubsection{Define mixture fraction. Which conditions have to be satisfied in order that mixture fraction becomes a useful variable? Explain the idea that, using mixture fraction, a non-premixed flame problem can be split into a mixing problem and a chemical problem.}

\subsubsection{Explain the structure of the strained laminar flamelet model for non-premixed combustion.}

\subsubsection{Explain how the influence of turbulent fluctuations of mixture fraction can be taken into account by making use of an assumed shape of its probability density function.}

\subsubsection{Explain why one is interested in calculating the variance of mixture fraction. How does the scalar dissipation rate influence the variance of mixture fraction?}

\subsubsection{Explain how model development and experiments come together in the process of model validation. Provide some conditions for good model validation.}

\subsubsection{List some key characteristics of the Sandia Flame series (D-E-F).}

\subsubsection{Discuss the presence and importance of stretch in the counterflow flame configuration.}

\subsubsection{How can a laminar premixed flame be described in an intrinsic way using a so-called progress variable? How can a turbulent premixed flame be described using a progress variable?}

\subsubsection{Provide an outline of the FGM approach to turbulent non-premixed combustion.}

\subsubsection{In which sense is the Transported PDF method to turbulent combustion completely different from flamelet models?}

\subsubsection{In which cases are flamelet models recommended, and in which cases transported PDF models?}

\subsection{Combustion Experiments and Diagnostics}

\subsubsection{What roles do canonical flames play in combustion research?}

\subsubsection{Laminar canonical flames are provided on standardized burners, for instance: the laminar premixed flame (Bunsen burner), the laminar ``rich'' premixed flame (McKenna burner), and the laminar diffusion flame (Wolfhard--Parker burner). Describe their laboratory flame geometries and the working principle of these burners. How can these burners be used to study combustion?}

\subsubsection{What are the three main advantages of optical diagnostics in flames?}

\subsubsection{What is temporal-, spatial-, and spectral resolution/control?}

\subsubsection{What is the Boltzmann thermal population distribution and how can it be used for thermometry in combustion?}

\subsubsection{How can species be detected and identified in flames? Why do we need this?}

\subsubsection{Give three examples of light--matter interaction and describe their working principles.}

\subsubsection{If you were given the assignment of measuring the laminar burning velocity of the simplest hydrocarbon fuel, methane, how would you do it?}

\subsection{Combustion Emissions}

\subsubsection{What are the main emissions from an aircraft engine?}

\subsubsection{How and why are they different from the emissions of cars and other terrestrial sources?}

\subsubsection{How do these emissions affect global warming?}

\subsubsection{What are the emission characteristics of a conventional aero engine combustor?}

\subsubsection{What are the different \ce{NOx} formation mechanisms?}

\subsubsection{What is the Zel'dovich mechanism for thermal \ce{NOx} formation? Discuss the importance of temperature fluctuations in thermal \ce{NOx} formation.}

\subsection{Practical Combustion Systems}

\subsubsection{What are the main methods of reducing \ce{NOx} emissions?}

\subsubsection{What is soot and what is the mechanism of soot formation in combustion?}

\subsubsection{What is soot growth and agglomeration?}

\subsubsection{What are the main design criteria and constraints in designing an aero engine combustor?}

\subsubsection{Explain the working principle of a conventional combustor, RQL combustor, staged combustor, and lean combustor.}

\subsection{Knowledge Integration and Application}

\subsubsection{Is it possible to design a burner that has a stable flame both when methane is used as a fuel and when hydrogen is used as fuel? What would be your suggestions to realize this objective and why?}

\subsubsection{Radiation plays a different role in the combustion process in rockets, gas turbines, and process furnaces. Explain why. How can we influence the type and amount of radiation produced by a flame?}

\subsubsection{List possible sources of pollution from natural gas combustion.}

\subsubsection{Explain how \ce{NOx} formation via the Zel'dovich mechanism can be described.}

\subsubsection{Explain how the burner design can influence thermal \ce{NOx} formation.}

\subsubsection{Explain how soot formation arises and how it could be modelled. Explain how the burner design can influence soot formation.}
